\chapter{Tầng 1: Thu thập và Chuẩn hóa Dữ liệu}
\section{Tổng quan Tầng Thu thập Dữ liệu}
Tầng 1 là lớp nền tảng của hệ thống STWI, có nhiệm vụ tiếp nhận dữ liệu thô từ phần cứng ngoại vi và chuyển hóa thành Tensor chuẩn hóa $\mathbf{X} \in \mathbb{R}^{B \times 12 \times N \times 16}$ để cung cấp cho Tầng 2. Dữ liệu đến từ hai nguồn chính: (1) Mạng lưới camera CCTV và (2) Hệ thống cảm biến môi trường và khí tượng.
\section{Phân hệ Thị giác Máy tính (CCTV Pipeline)}
\subsection{Kiến trúc Pipeline CCTV}
\begin{figure}[H]
\centering
\begin{tikzpicture}[node distance=0.5cm and 0.5cm, scale=0.72, transform shape]
  \node[sysbox=tier1col, text width=2.4cm, minimum height=1.5cm, font=\scriptsize] (cam)
    {Camera demo $\leq20$ luồng\\Ghi sẵn/RTSP\\Không lưu video thô};
  \node[sysbox=tier2col, right=0.7cm of cam, text width=2.4cm, minimum height=1.5cm, font=\scriptsize] (yolo)
    {YOLOv8 local\\Object Detection\\Private weights};
  \node[sysbox=tier2col, right=0.7cm of yolo, text width=2.4cm, minimum height=1.5cm, font=\scriptsize] (bt)
    {ByteTrack\\Multi-Object\\Tracking};
  \node[sysbox=tier3col, right=0.7cm of bt, text width=2.6cm, minimum height=1.5cm, font=\scriptsize] (agg)
    {Aggregation\\Chu kỳ 5 phút\\/camera};
  \node[sysbox=tier4col, right=0.7cm of agg, text width=2.4cm, minimum height=1.5cm, font=\scriptsize] (tsdb)
    {TimescaleDB\\Aggregate 5 phút\\Runtime};
  \draw[arrow] (cam) -- (yolo);
  \draw[arrow] (yolo) -- (bt);
  \draw[arrow] (bt) -- (agg);
  \draw[arrow] (agg) -- (tsdb);
  % Output labels
  \node[below=0.6cm of agg, font=\tiny\itshape, text=codegray, text width=2.8cm, align=center]
    {3 biến:\\traffic\_volume\_5m\\avg\_speed\_kmh\\heavy\_vehicle\_ratio};
  \node[below=0.6cm of tsdb, font=\tiny\itshape, text=codegray, text width=2.4cm, align=center]
    {JSON records\\timestamped};
\end{tikzpicture}
\caption{Luồng xử lý dữ liệu CCTV --- từ Camera đến TimeSeries DB}
\label{fig:cctv_pipeline}
\end{figure}
\subsection{YOLOv8 -- Phát hiện Phương tiện}
YOLOv8 \cite{jocher2023yolo} là kiến trúc phát hiện vật thể (object detection) thời gian thực tiên tiến, xử lý mỗi khung hình (frame) trong một lần forward pass duy nhất:
\begin{itemize}[leftmargin=1.5cm, itemsep=4pt]
  \item \textbf{Backbone:} CSPDarknet53 với C2f modules cho trích xuất đặc trưng đa tỉ lệ.
  \item \textbf{Neck:} PANet (Path Aggregation Network) để kết hợp đặc trưng từ nhiều tỉ lệ.
  \item \textbf{Head:} Decoupled head dự báo bounding box và class score độc lập.
  \item \textbf{Tốc độ inference:} đo theo benchmark nội bộ của artifact được promote; không dùng số liệu khi thiếu phần cứng, batch và phiên bản model.
\end{itemize}
Trong MVP, YOLOv8 được train local từ dataset private theo runbook vision nội bộ. Roboflow export chỉ là nguồn dữ liệu offline; runtime chính dùng weight local có version, checksum, class map, metrics và privacy review. Hosted workflow inference của Roboflow chỉ là adapter tùy chọn, không phải phụ thuộc triển khai.

Các lớp phương tiện cần nhận diện:
\begin{table}[H]
\centering
\caption{Phân loại phương tiện trong STWI CCTV Pipeline}
\label{tab:vehicle_class}
\renewcommand{\arraystretch}{1.4}
\begin{tabular}{>{\centering}p{1.5cm} p{3.5cm} p{5cm} p{2cm}}
\toprule
\textbf{Lớp ID} & \textbf{Loại phương tiện} & \textbf{Ví dụ cụ thể} & \textbf{PCU$^*$} \\
\midrule
0 & Ô tô con & Sedan, SUV, MPV & 1.0 \\
\rowcolor{rowgray}
1 & Xe máy & Xe ga, xe số, xe đạp điện & 0.5 \\
2 & Xe buýt & Buýt thường, buýt nhanh BRT & 2.5 \\
\rowcolor{rowgray}
3 & Xe tải & Tải nhỏ, tải trung, container & 2.0--3.5 \\
\bottomrule
\multicolumn{4}{l}{\small $^*$PCU: Passenger Car Unit --- đơn vị quy đổi ra xe con tương đương}
\end{tabular}
\end{table}
\subsection{ByteTrack -- Theo dõi Đa đối tượng}
ByteTrack \cite{zhang2020bytetrack} giải quyết vấn đề đếm trùng phương tiện giữa các frame liên tiếp bằng cách theo dõi \textit{tất cả} bounding box (bao gồm cả các detection có confidence thấp):
\begin{enumerate}[leftmargin=1.5cm, itemsep=4pt]
  \item \textbf{High-confidence detections} ($\text{conf} \geq \theta_{\text{high}}$): Ghép với tracks hiện có theo IoU (Hungarian algorithm).
  \item \textbf{Low-confidence detections} ($\theta_{\text{low}} \leq \text{conf} < \theta_{\text{high}}$): Thử ghép với các tracks bị miss ở bước trên.
  \item \textbf{Kalman Filter}: Dự báo vị trí track trong frame tiếp theo để xử lý occlusion.
\end{enumerate}
Kết quả: Mỗi phương tiện được gán ID duy nhất, cho phép tính lưu lượng thực (không đếm trùng) và vận tốc trung bình từ quỹ đạo di chuyển.
\subsection{Các Biến Đầu ra từ CCTV}
Sau mỗi chu kỳ 5 phút, pipeline CCTV xuất 3 biến cho mỗi nút giao:
\begin{table}[H]
\centering
\caption{Biến đặc trưng từ CCTV Pipeline}
\label{tab:cctv_features}
\renewcommand{\arraystretch}{1.4}
\begin{tabular}{p{4.2cm} p{2.2cm} p{5.6cm}}
\toprule
\textbf{Tên biến} & \textbf{Đơn vị} & \textbf{Mô tả} \\
\midrule
\texttt{traffic\_volume\_5m} & vehicles/5min & Tổng số phương tiện đi qua nút giao trong cửa sổ 5 phút \\
\rowcolor{rowgray}
\texttt{avg\_speed\_kmh} & km/h & Tốc độ trung bình, chỉ phát hành khi camera có homography hoặc hệ số hiệu chỉnh được duyệt \\
\texttt{heavy\_vehicle\_ratio} & $[0,1]$ & Tỷ lệ xe tải và xe buýt trong tổng lưu lượng \\
\bottomrule
\end{tabular}
\end{table}
\section{Phân hệ Cảm biến Môi trường và Khí tượng}
\subsection{Mạng lưới Cảm biến}
Hệ thống tích hợp dữ liệu từ hai loại cảm biến:
\begin{table}[H]
\centering
\caption{Danh sách biến cảm biến môi trường và khí tượng}
\label{tab:sensors}
\renewcommand{\arraystretch}{1.4}
\begin{tabular}{p{2cm} p{3.2cm} p{2.5cm} p{4cm}}
\toprule
\textbf{Nhóm} & \textbf{Biến} & \textbf{Ký hiệu} & \textbf{Đơn vị / Dải giá trị} \\
\midrule
\multirow{5}{*}{Khí thải} & Carbon monoxide & CO & $\mu$g/m$^3$, 0--500 \\
  & Carbon dioxide & CO$_2$ & ppm, 300--5.000 \\
  & Nitrogen oxides & NO$_x$ & $\mu$g/m$^3$, 0--200 \\
  & Bụi mịn nhỏ & PM$_{2.5}$ & $\mu$g/m$^3$, 0--250 \\
  & Bụi mịn lớn & PM$_{10}$ & $\mu$g/m$^3$, 0--400 \\
\midrule
\rowcolor{rowgray}
\multirow{3}{*}{Khí tượng} & Nhiệt độ & $T$ & $^\circ$C, 15--42 \\
\rowcolor{rowgray}
  & Độ ẩm & $H$ & \%, 30--100 \\
\rowcolor{rowgray}
  & Tốc độ gió & $W$ & m/s, 0--20 \\
\bottomrule
\end{tabular}
\end{table}
\begin{tcolorbox}[tipbox]
\textbf{Logic kỹ thuật quan trọng:} Tốc độ gió $W$ có mối quan hệ phi tuyến với PM$_{2.5}$ và PM$_{10}$: gió mạnh ($W > 5$ m/s) phát tán bụi nhanh, nhưng gió yếu ($W < 1$ m/s) cho phép bụi tích tụ. Mô hình GCN--LSTM cần nắm bắt được mối quan hệ này để dự báo chất lượng không khí trong điều kiện ùn tắc kéo dài.
\end{tcolorbox}
\subsection{Giao thức Thu thập và Xử lý Mất dữ liệu}
Cảm biến giao tiếp qua MQTT với broker trung tâm. Chiến lược xử lý dữ liệu thiếu:
\begin{itemize}[leftmargin=1.5cm, itemsep=3pt]
  \item \textbf{Ngắt $< 15$ phút:} Linear interpolation giữa hai điểm đo lân cận.
  \item \textbf{Ngắt 15--60 phút:} Sử dụng giá trị trung bình theo giờ tương ứng từ lịch sử 30 ngày.
  \item \textbf{Ngắt $> 60$ phút:} Đánh dấu \texttt{sensor\_unavailable}, giữ node order, điền giá trị theo policy và đặt $M=0$; quality thấp chuyển degraded/\texttt{needs\_review}.
\end{itemize}
\section{Cấu trúc và Chuẩn hóa Dữ liệu}
\subsection{16 Biến Đặc trưng (Feature Engineering)}
Sau khi thu thập, dữ liệu được ghép thành vector đặc trưng 16 chiều cho mỗi nút giao tại mỗi bước thời gian 5 phút:
\begin{table}[H]
\centering
\caption{16 biến đặc trưng đầu vào cho mô hình GCN–LSTM}
\label{tab:features}
\renewcommand{\arraystretch}{1.4}
\begin{tabular}{>{\centering}p{1.5cm} >{\centering}p{0.8cm} p{3.5cm} p{5cm}}
\toprule
\textbf{Nhóm} & \textbf{ID} & \textbf{Tên biến} & \textbf{Mô tả} \\
\midrule
\multirow{3}{*}{Giao thông} & F1 & \texttt{traffic\_volume\_5m} & Lưu lượng xe (vehicles/5 min) \\
  & F2 & \texttt{avg\_speed\_kmh} & Vận tốc trung bình (km/h) \\
  & F3 & \texttt{heavy\_vehicle\_ratio} & Tỷ lệ xe tải/nặng $(0, 1)$ \\
\midrule
\rowcolor{rowgray}
\multirow{5}{*}{Khí thải} & F4 & \texttt{co\_ppm} & Nồng độ CO (ppm) \\
\rowcolor{rowgray}
  & F5 & \texttt{co2\_ppm} & Nồng độ CO$_2$ (ppm) \\
\rowcolor{rowgray}
  & F6 & \texttt{nox\_ppb} & Nồng độ NO$_x$ (ppb) \\
\rowcolor{rowgray}
  & F7 & \texttt{pm25\_ugm3} & Bụi mịn PM$_{2.5}$ ($\mu$g/m$^3$) \\
\rowcolor{rowgray}
  & F8 & \texttt{pm10\_ugm3} & Bụi mịn PM$_{10}$ ($\mu$g/m$^3$) \\
\midrule
\multirow{3}{*}{Khí tượng} & F9 & \texttt{temperature\_c} & Nhiệt độ ($^\circ$C) \\
  & F10 & \texttt{humidity\_pct} & Độ ẩm (\%) \\
  & F11 & \texttt{wind\_speed\_ms} & Tốc độ gió (m/s) \\
\midrule
\rowcolor{rowgray}
\multirow{5}{*}{Thời gian} & F12 & \texttt{time\_of\_day\_sin} & $\sin(2\pi \cdot \text{minute}/1440)$ \\
  & F13 & \texttt{time\_of\_day\_cos} & $\cos(2\pi \cdot \text{minute}/1440)$ \\
  & F14 & \texttt{day\_of\_week\_sin} & $\sin(2\pi \cdot \text{day}/7)$ \\
  & F15 & \texttt{day\_of\_week\_cos} & $\cos(2\pi \cdot \text{day}/7)$ \\
\midrule
\rowcolor{rowgray}
Tín hiệu & F16 & \texttt{green\_time\_ratio} & Tỷ lệ thời gian đèn xanh trong cửa sổ 5 phút \\
\bottomrule
\end{tabular}
\end{table}
\subsection{Chuẩn hóa và Quality Mask}
Scaler chỉ được fit trên training split theo thời gian và chỉ áp dụng cho các feature liên tục. Các biến cyclical và ratio giữ nguyên miền giá trị. Mọi giá trị thiếu hoặc được impute có mask bằng 0.
\begin{tcolorbox}[dangerbox]
\textbf{Hợp đồng dữ liệu:}
$\mathbf{X}\in\mathbb{R}^{B\times12\times N\times16}$,
$\mathbf{M}\in\{0,1\}^{B\times12\times N\times16}$ và
$\mathbf{A}\in\mathbb{R}^{N\times N}$ phải dùng cùng node order trong \texttt{node\_registry}.
\end{tcolorbox}
Feature thứ 16 là \texttt{green\_time\_ratio}, tức tỷ lệ thời gian đèn xanh trong aggregate 5 phút; không dùng trạng thái đèn tức thời.
\subsection{Xây dựng Tensor 4 chiều}
\begin{lstlisting}[caption={Xây dựng tensor mạng và missing mask},label={lst:tensor_build}]
import torch
def build_network_batch(windows, masks, adjacency):
    # windows/masks: [B, 12, N, 16]; adjacency: [N, N]
    X = torch.as_tensor(windows, dtype=torch.float32)
    M = torch.as_tensor(masks, dtype=torch.bool)
    A = torch.as_tensor(adjacency, dtype=torch.float32)
    assert X.ndim == 4 and X.shape[1] == 12 and X.shape[3] == 16
    assert M.shape == X.shape
    assert A.shape == (X.shape[2], X.shape[2])
    return X, M, A
\end{lstlisting}
Batch là tập cửa sổ của \emph{toàn mạng}; trục $N$ mới là node/camera đã ánh xạ. Dữ liệu thiếu ngắn hạn được impute nhưng luôn giữ $\mathbf{M}=0$ để model và metrics nhận biết chất lượng.
\section{Kiến trúc Lưu trữ và Đồng bộ Dữ liệu}
Để đảm bảo đồng bộ dữ liệu từ nhiều nguồn (CCTV và cảm biến có thể có độ trễ khác nhau), pipeline sử dụng kiến trúc \textbf{Event-Driven + Time-Window Buffering}:
\begin{enumerate}[leftmargin=1.5cm, itemsep=4pt]
  \item Tất cả messages được publish lên message broker (MQTT) với timestamp chính xác.
  \item Time-window buffer (5 phút) tích lũy dữ liệu từ tất cả sensors tại một nút.
  \item Sau khi cửa sổ đóng, data aggregator tính toán trung bình và ghép thành vector 16 chiều.
  \item Nếu một cảm biến thiếu dữ liệu, áp dụng chiến lược imputation tương ứng.
  \item Vector hoàn chỉnh được ghi vào TimescaleDB với key `\texttt{node\_id + timestamp}'.
\end{enumerate}
